%=====================================================================
\chapter{Linear Regression}\label{ch:linreg}
%=====================================================================
\locator{2}{Part II --- Linear Models and Generalization}

\begin{outcomes}
By the end of this chapter you will be able to:
\begin{tight}
  \item \textbf{Write} the linear regression model and its mean squared error
        loss, and \textbf{interpret} each weight.
  \item \textbf{Fit} a simple linear regression by hand with the least-squares
        formulas, and \textbf{state} the normal equation for many features.
  \item \textbf{Derive} the gradient descent update for linear regression and
        \textbf{perform} a step by hand.
  \item \textbf{Extend} the model with multiple and polynomial features.
  \item \textbf{Evaluate} a regression model with MAE, RMSE and $R^2$ against a
        baseline, and \textbf{read} a residual plot.
\end{tight}
\end{outcomes}

\begin{prereq}
\chref{math} (dot products, derivatives, gradient descent, maximum likelihood)
and \chref{workflow} (splits, scaling, baselines).
\end{prereq}

\begin{keyterms}
regression $\bullet$ linear model $\bullet$ weight $\bullet$ bias (intercept)
$\bullet$ residual $\bullet$ mean squared error $\bullet$ least squares
$\bullet$ normal equation $\bullet$ batch, stochastic and mini-batch gradient
descent $\bullet$ multiple regression $\bullet$ polynomial features $\bullet$
MAE $\bullet$ RMSE $\bullet$ $R^2$ $\bullet$ residual plot
\end{keyterms}

\section{The Model}

Linear regression predicts a number as a weighted sum of the features plus a
constant:
\[
\hat{y} = w_1 x_1 + w_2 x_2 + \dots + w_d x_d + b = \mathbf{w}\cdot\mathbf{x} + b
\]
With one feature this is the straight line $\hat{y} = wx + b$ of school algebra.
The \term{weight} $w_j$ is the change in the prediction for a one-unit increase in
feature $j$ with the other features held fixed; the \term{bias} (or intercept)
$b$ is the prediction when every feature is zero.

Linear regression is the first model in the course because every later model
either extends it (logistic regression, the perceptron, the SVM, a neural
network's output layer) or is compared against it. It is also, very often, the
right answer: fast, stable, and explainable in one sentence per feature.

\section{The Loss: Mean Squared Error}

\begin{definitionbox}[Residual and Mean Squared Error]
The \term{residual} of example $i$ is $e_i = y_i - \hat{y}_i$, the vertical gap
between the observed value and the prediction. The \term{mean squared error} of
the model on $n$ examples is
\[
\text{MSE}(\mathbf{w}, b) = \frac{1}{n}\sum_{i=1}^{n}\left(y_i - \hat{y}_i\right)^2
\]
Fitting a linear regression means choosing $\mathbf{w}$ and $b$ to make the MSE
as small as possible --- the method of \term{least squares}.
\end{definitionbox}

\dsfig{%
\begin{tikzpicture}
\begin{axis}[width=8.6cm, height=5.6cm, axis lines=left,
  xlabel={story points $x$}, ylabel={hours $y$},
  tick label style={font=\tiny}, label style={font=\scriptsize},
  xmin=0, xmax=6, ymin=0, ymax=14, xtick={1,...,5}]
\addplot[primary, line width=1.3pt, domain=0.3:5.7] {2.2*x+0.4};
\foreach \xx/\yy in {1/3,2/5,3/6,4/9,5/12}{
  \edef\tmp{\noexpand\draw[accent, line width=0.9pt]
    (axis cs:\xx,\yy) -- (axis cs:\xx,{2.2*\xx+0.4});}\tmp}
\addplot[only marks, mark=*, mark size=2.2pt, secondary] coordinates
  {(1,3)(2,5)(3,6)(4,9)(5,12)};
\node[font=\tiny, primary, anchor=west] at (axis cs:4.2,12.6) {$\hat{y} = 2.2x + 0.4$};
\node[font=\tiny, accent, anchor=west] at (axis cs:3.1,5.6) {residual $e_3=-1.0$};
\end{axis}
\end{tikzpicture}}%
{The least-squares line through five user stories. Each red segment is a residual; the
line is the one that makes the sum of their squares smallest.}{least-squares}

\begin{conceptbox}[Why Square the Residuals?]
Three reasons. Squaring makes every error positive, so over- and under-predictions
cannot cancel. It makes the loss smooth, so it has a derivative everywhere and a
single minimum that a formula can find. And it is what maximum likelihood gives
(\chref{math}): if the noise around the true line is Gaussian, the most likely
line is exactly the least-squares line. The price is that one large error
counts a great deal --- an error of 10 costs a hundred times an error of 1 ---
so least squares is pulled hard by outliers.
\end{conceptbox}

\section{Least Squares by Formula}

For one feature, setting the derivatives of the MSE with respect to $w$ and $b$
to zero gives a closed-form answer:
\[
w = \frac{\sum_i (x_i - \bar{x})(y_i - \bar{y})}{\sum_i (x_i - \bar{x})^2},
\qquad b = \bar{y} - w\bar{x}
\]
The line always passes through the point of means $(\bar{x}, \bar{y})$.

\begin{worked}[Fitting Hours to Story Points]
Five completed user stories: story points $x = 1, 2, 3, 4, 5$ took
$y = 3, 5, 6, 9, 12$ hours.

\smallskip
\textbf{Means:} $\bar{x} = 3$, $\bar{y} = 35/5 = 7$.

\textbf{Deviations:} $x - \bar{x} = -2, -1, 0, 1, 2$ and $y - \bar{y} = -4, -2, -1,
2, 5$.

\textbf{Slope:}
\[
w = \frac{(-2)(-4) + (-1)(-2) + (0)(-1) + (1)(2) + (2)(5)}{4 + 1 + 0 + 1 + 4}
  = \frac{8 + 2 + 0 + 2 + 10}{10} = \mathbf{2.2}
\]
\textbf{Intercept:} $b = 7 - 2.2 \times 3 = \mathbf{0.4}$.

\smallskip
\textbf{Reading it:} each extra story point adds 2.2 hours; the intercept, 0.4
hours, is the fixed overhead of any story. A new 6-point story is predicted to
take $2.2 \times 6 + 0.4 = 13.6$ hours.

\textbf{Residuals:} predictions $2.6, 4.8, 7.0, 9.2, 11.4$ give residuals
$0.4, 0.2, -1.0, -0.2, 0.6$, which sum to zero --- they always do for a
least-squares line with an intercept. The sum of squares is
$0.16 + 0.04 + 1 + 0.04 + 0.36 = 1.6$, so $\text{MSE} = 0.32$.
\end{worked}

\begin{definitionbox}[The Normal Equation]
With $d$ features, put a column of ones in front of $X$ to absorb the bias, so the
model is $\hat{\mathbf{y}} = X\boldsymbol{\theta}$. The least-squares parameters
solve the \term{normal equation}
\[
X^{\top}X\,\boldsymbol{\theta} = X^{\top}\mathbf{y}
\quad\Longrightarrow\quad
\boldsymbol{\theta} = (X^{\top}X)^{-1}X^{\top}\mathbf{y}
\]
For the five stories,
\[
X^{\top}X = \begin{pmatrix} 5 & 15 \\ 15 & 55\end{pmatrix}, \qquad
X^{\top}\mathbf{y} = \begin{pmatrix} 35 \\ 127 \end{pmatrix},
\]
whose solution is $b = 0.4$, $w = 2.2$ as before.
\end{definitionbox}

The normal equation is exact and needs no learning rate, but inverting a
$d \times d$ matrix costs roughly $d^3$ operations, and it fails outright when
two features are perfect copies of each other ($X^{\top}X$ has no inverse). For
many features, or data that does not fit in memory, gradient descent is used
instead.

\section{Fitting by Gradient Descent}

Differentiating the MSE with the chain rule gives the gradient:
\[
\frac{\partial\,\text{MSE}}{\partial w_j} = \frac{2}{n}\sum_{i=1}^{n}(\hat{y}_i - y_i)\,x_{ij},
\qquad
\frac{\partial\,\text{MSE}}{\partial b} = \frac{2}{n}\sum_{i=1}^{n}(\hat{y}_i - y_i)
\]
Read it in words: each weight moves in proportion to the average of
\emph{(error) $\times$ (its feature)}. A feature that is large exactly when the
model under-predicts has its weight pushed up.

\begin{worked}[One Gradient Descent Step]
Start from $w = 0$, $b = 0$ on the five stories, with $\alpha = 0.05$. Every
prediction is 0, so the errors $\hat{y} - y$ are $-3, -5, -6, -9, -12$ and the
MSE is $(9 + 25 + 36 + 81 + 144)/5 = 59$.
\[
\frac{\partial}{\partial w} = \frac{2}{5}\big[(-3)(1) + (-5)(2) + (-6)(3) + (-9)(4) + (-12)(5)\big]
= \frac{2}{5}(-127) = -50.8
\]
\[
\frac{\partial}{\partial b} = \frac{2}{5}(-3 - 5 - 6 - 9 - 12) = \frac{2}{5}(-35) = -14
\]
\textbf{Update:} $w \leftarrow 0 - 0.05(-50.8) = 2.54$, $\;b \leftarrow 0 - 0.05(-14)
= 0.70$.

The new predictions $3.24, 5.78, 8.32, 10.86, 13.40$ give an MSE of $2.29$ --- down
from 59 in one step. Repeating the update converges to $w = 2.2$, $b = 0.4$, the
same answer as the formula.
\end{worked}

\Needspace{9\baselineskip}
\rowcolors{2}{white}{lightbg}
\begin{longtable}{@{}L{3.3cm}L{5.4cm}L{6.3cm}@{}}
\toprule
\rowcolor{primary!15}
\textbf{Variant} & \textbf{Each step uses} & \textbf{Trade-off} \\
\midrule
\endhead
Batch & All $n$ examples & Smooth, exact steps; slow when $n$ is large \\
Stochastic (SGD) & One example, chosen at random & Very cheap steps; noisy path
that jitters around the minimum \\
Mini-batch & A batch of, say, 32--256 examples & The practical default; what
neural networks use (\chref{ann}) \\
\bottomrule
\end{longtable}

\begin{alertbox}[Scale the Features Before Gradient Descent]
If one feature is measured in thousands and another in fractions, the MSE surface
is a long narrow valley: a learning rate small enough for the steep direction
crawls along the shallow one. Standardising the features (\chref{workflow})
makes the valley round and gradient descent converge in far fewer steps. The
normal equation does not care.
\end{alertbox}

\section{Multiple and Polynomial Features}

With several features each weight is a \emph{partial} effect: the change in
$\hat{y}$ for one more unit of $x_j$ \emph{with the other features held fixed}.
That is why the weight of a feature can change sign when another, correlated
feature is added.

A linear model can also fit curves. Add powers of a feature as new features ---
$x$, $x^2$, $x^3$ --- and fit
$\hat{y} = w_1 x + w_2 x^2 + w_3 x^3 + b$. The model is still \emph{linear in the
weights}, so everything in this chapter applies unchanged; only the features are
new. The danger is obvious as soon as it is said: with enough powers, the curve
can pass through every training point. \chref{generalization} begins exactly there.

\section{Measuring a Regression Model}

\begin{definitionbox}[MAE, RMSE and $R^2$]
\term{Mean absolute error}: $\text{MAE} = \frac{1}{n}\sum|y_i - \hat{y}_i|$ --- the
typical size of an error, in the units of $y$.\\[3pt]
\term{Root mean squared error}: $\text{RMSE} = \sqrt{\text{MSE}}$ --- also in the
units of $y$, but weighting large errors more.\\[3pt]
\term{Coefficient of determination}:
$R^2 = 1 - \dfrac{\sum(y_i - \hat{y}_i)^2}{\sum(y_i - \bar{y})^2}$ --- the fraction
of the variance in $y$ the model explains. 1 is perfect; 0 is no better than
always predicting the mean; on test data it can be negative.
\end{definitionbox}

\begin{worked}[Scoring the Story-Point Model]
Residuals $0.4, 0.2, -1.0, -0.2, 0.6$.
\begin{tight}
  \item $\text{MAE} = (0.4 + 0.2 + 1.0 + 0.2 + 0.6)/5 = \mathbf{0.48}$ hours.
  \item $\text{RMSE} = \sqrt{1.6/5} = \sqrt{0.32} = \mathbf{0.57}$ hours.
  \item $\sum(y - \bar{y})^2 = 16 + 4 + 1 + 4 + 25 = 50$, so
        $R^2 = 1 - 1.6/50 = \mathbf{0.968}$.
\end{tight}
\smallskip
\textbf{Caution.} These are \emph{training} scores on five points fitted by two
parameters. They say the line describes these five stories well. Whether it
predicts the next story well can only be answered on stories it has not seen.
\end{worked}

\dsfig{%
\begin{tikzpicture}
\begin{groupplot}[
  group style={group size=2 by 1, horizontal sep=1.3cm},
  width=6.7cm, height=4.6cm, axis lines=left,
  tick label style={font=\tiny}, label style={font=\tiny},
  title style={font=\scriptsize\bfseries},
  xlabel={prediction $\hat{y}$}, ylabel={residual $y-\hat{y}$},
  xmin=0, xmax=10, ymin=-3, ymax=3]
\nextgroupplot[title={Healthy: no pattern}]
\addplot[gray!50, dashed] coordinates {(0,0)(10,0)};
\addplot[only marks, mark=*, mark size=1.4pt, greenhl!70!black] coordinates {
 (0.6,0.8)(1.1,-0.9)(1.7,0.3)(2.2,-0.4)(2.9,1.1)(3.3,-1.2)(3.8,0.2)(4.4,0.9)
 (4.9,-0.6)(5.3,0.4)(5.9,-1.0)(6.4,1.3)(6.8,-0.2)(7.3,0.6)(7.9,-0.8)(8.4,0.1)
 (8.8,1.0)(9.3,-0.5)};
\nextgroupplot[title={A curve the line missed}]
\addplot[gray!50, dashed] coordinates {(0,0)(10,0)};
\addplot[only marks, mark=*, mark size=1.4pt, accent] coordinates {
 (0.6,2.2)(1.1,1.6)(1.7,1.2)(2.2,0.4)(2.9,0.1)(3.3,-0.6)(3.8,-0.8)(4.4,-1.3)
 (4.9,-1.1)(5.3,-1.4)(5.9,-1.0)(6.4,-0.9)(6.8,-0.5)(7.3,-0.1)(7.9,0.6)(8.4,1.0)
 (8.8,1.7)(9.3,2.3)};
\addplot[accent!50, line width=0.9pt, domain=0.4:9.5, samples=40] {0.105*(x-5)^2-1.35};
\end{groupplot}
\end{tikzpicture}}%
{Residual plots. Left: residuals scattered evenly around zero, as they should be. Right:
a U-shape, meaning the true relationship is curved; adding an $x^2$ feature would remove
it.}{residual-plots}

\begin{alertbox}[Always Compare With the Mean]
The baseline for regression is ``predict the training mean for everyone'', which
has $R^2 = 0$. On the 400 completed stories of this chapter's lab it has an RMSE
of 11.8 hours; linear regression brings that to 4.3. Whether that improvement is worth having is
a question about the application, not the algorithm --- but without the baseline
nobody could even ask it.
\end{alertbox}

\begin{pitfall}
\begin{tight}
  \item \textbf{Reading a weight as a cause.} $w = 2.2$ hours per point describes
        the data; it does not say that adding points to a story would add hours.
        \chref{ethics} returns to this.
  \item \textbf{Comparing raw weights across features.} A weight's size depends
        on its feature's units. Standardise first, then compare.
  \item \textbf{Extrapolating.} The line was fitted on 1--5 points; a 40-point
        epic is outside anything it has seen.
  \item \textbf{Ignoring the residual plot.} A good $R^2$ with a curved residual
        pattern means a better model is available.
  \item \textbf{Letting one outlier drive the fit.} Squared loss gives a single
        wild point enormous influence. Look at the data first.
\end{tight}
\end{pitfall}

%---------------------------------------------------------------------
\begin{lab}[Least Squares Three Ways]
The worked example by formula, then by gradient descent, then a larger dataset:
400 completed user stories, with the hours each took predicted from five
features --- story points, files touched, requirement changes, the developer's
years of experience, and whether the story touches the database.
\begin{lstlisting}[language=Python]
import numpy as np
from sklearn.model_selection import train_test_split
from sklearn.preprocessing import StandardScaler
from sklearn.linear_model import LinearRegression
from sklearn.metrics import mean_absolute_error, mean_squared_error, r2_score

# --- 1. the worked example, by formula -----------------------------------
x = np.array([1, 2, 3, 4, 5.0])          # story points
y = np.array([3, 5, 6, 9, 12.0])         # hours
w = ((x - x.mean()) * (y - y.mean())).sum() / ((x - x.mean()) ** 2).sum()
b = y.mean() - w * x.mean()
print(f"least squares:    w = {w:.2f}, b = {b:.2f}")

# --- 2. the same line by gradient descent --------------------------------
wg, bg, alpha = 0.0, 0.0, 0.05
for step in range(2000):
    err = (wg * x + bg) - y              # y_hat - y
    wg -= alpha * 2 * (err * x).mean()
    bg -= alpha * 2 * err.mean()
print(f"gradient descent: w = {wg:.2f}, b = {bg:.2f}")

# --- 3. 400 completed stories, against the mean baseline -----------------
rng = np.random.default_rng(5)
n = 400
points = rng.choice([1, 2, 3, 5, 8, 13], n)
files = rng.poisson(2 + 0.5 * points)            # files touched
changes = rng.poisson(1.0, n)                    # requirement changes
years = rng.integers(0, 11, n)                   # developer experience
db = (rng.random(n) < 0.3).astype(int)           # touches the database
t = (2 + 2.2 * points + 0.8 * files + 1.5 * changes - 0.3 * years
     + 3 * db + rng.normal(0, 4, n)).clip(min=1)  # actual hours
X = np.column_stack([points, files, changes, years, db])
names = ["points", "files", "changes", "years", "db"]
X_tr, X_te, t_tr, t_te = train_test_split(X, t, test_size=0.25,
                                          random_state=0)
scaler = StandardScaler().fit(X_tr)
model = LinearRegression().fit(scaler.transform(X_tr), t_tr)
pred = model.predict(scaler.transform(X_te))
base = np.full_like(t_te, t_tr.mean())   # always predict the mean
for name, p in [("baseline (mean)", base), ("linear regression", pred)]:
    print(f"{name:18s} MAE {mean_absolute_error(t_te, p):5.1f}  "
          f"RMSE {mean_squared_error(t_te, p) ** 0.5:5.1f}  "
          f"R2 {r2_score(t_te, p):5.2f}")
for i in np.argsort(-np.abs(model.coef_))[:3]:
    print(f"  weight on {names[i]:7s}: {float(model.coef_[i]):6.1f}")
\end{lstlisting}
Change \texttt{alpha} to \texttt{0.2} in part 2 and run it again. Explain what you
see using the worked example of \chref{math}.
\end{lab}

%---------------------------------------------------------------------
\begin{checkpoint}
\begin{tightnum}
  \item Fit $\hat{y} = wx + b$ to the points $(0, 1)$, $(1, 3)$, $(2, 5)$ by
        the least-squares formulas.
  \item A model has $R^2 = -0.2$ on the test set. What does that mean?
  \item Why is $\hat{y} = w_1 x + w_2 x^2 + b$ still called a \emph{linear} model?
\end{tightnum}
\end{checkpoint}

%---------------------------------------------------------------------
\begin{chaptersummary}
\begin{tight}
  \item Linear regression predicts $\hat{y} = \mathbf{w}\cdot\mathbf{x} + b$; each
        weight is the effect of one more unit of its feature, others held fixed.
  \item Least squares minimises the mean squared error; it is the maximum
        likelihood fit under Gaussian noise, and it is sensitive to outliers.
  \item One feature: $w = \sum(x-\bar{x})(y-\bar{y}) / \sum(x-\bar{x})^2$,
        $b = \bar{y} - w\bar{x}$. Many: the normal equation.
  \item Gradient descent moves each weight by the average of error times feature;
        scale features first. Mini-batch is the practical default.
  \item Polynomial features let a linear model fit curves --- and overfit.
  \item Report MAE or RMSE in the units of $y$, and $R^2$, beside the
        mean baseline; check the residual plot for patterns.
\end{tight}
\end{chaptersummary}

\begin{reviewq}
  \item \textbf{(MCQ)} The least-squares line always passes through:
        (a)~the origin (b)~$(\bar{x}, \bar{y})$ (c)~the first data point (d)~the
        largest residual.
  \item \textbf{(MCQ)} Which metric is most affected by one very large error?
        (a)~MAE (b)~RMSE (c)~the median absolute error (d)~all equally.
  \item \textbf{(MCQ)} $R^2 = 0$ means the model: (a)~is perfect (b)~does no better
        than predicting the mean (c)~has no bias (d)~is overfitted.
  \item \textbf{(Short)} Explain in words the gradient of the MSE with respect to
        one weight.
  \item \textbf{(Short)} Give two situations where you would use gradient descent
        rather than the normal equation.
  \item \textbf{(Short)} Why can the weight of a feature change sign when another
        feature is added to the model?
  \item \textbf{(Applied)} For the points $(1, 2)$, $(2, 2)$, $(3, 5)$, $(4, 6)$,
        compute the least-squares line, the residuals, the MAE, the RMSE and $R^2$.
  \item \textbf{(Applied)} Starting from $w = 0$, $b = 0$ with $\alpha = 0.1$,
        perform one gradient descent step on the data of the previous question
        and compute the MSE before and after.
\end{reviewq}
